\chapter{Energy in Africa}


\section{Introduction: Africa and the Global Energy Context}

\subsection{Africa as a Continent}

Africa is frequently misrepresented in scale. To appreciate the challenges it faces, one must first grasp its sheer geographic magnitude. The distance from Egypt to South Africa spans approximately 6{,}435~km, from Algeria to South Africa around 6{,}894~km, and from Gambia to Somalia roughly 6{,}796~km — all figures that dwarf the Portugal-to-Moscow distance of only 3{,}768~km. Africa is thus a continent of vast internal distances, hosting 54 countries with deeply heterogeneous political, economic, and energy landscapes.

In 2025, the world population reached approximately 8.2~billion people. Africa accounted for 19\% of this total, second only to Asia (58\%), and ahead of the Americas (13\%), Europe (9\%), and Oceania (1\%). This demographic weight, combined with rapid population growth, makes the continent's energy future a global concern.

\subsection{Global Primary Energy Resources}

At the global scale, the energy mix has remained heavily dominated by fossil fuels (coal, oil, and natural gas) over the past five decades. Data spanning from 1970 to 2022 consistently show fossil fuels accounting for roughly 80 to 95\% of primary energy supply, with renewables growing modestly and nuclear energy contributing a relatively stable but small share. Despite accelerating renewable deployment in recent years, fossils remain the backbone of global energy production.

In 2020, global primary energy consumption stood at 556.6~EJ, with oil at 31\%, natural gas at 25\%, coal at 27\%, hydroelectricity at 7\%, nuclear at 4\%, and renewables at 6\%. That same year, global electricity production totalled 26{,}823~TWh, with coal leading at 35\%, followed by natural gas (23\%), hydroelectricity (16\%), nuclear (10\%), renewables (12\%), and oil (3\%).

\subsection{Energy Conversion Losses}

A striking illustration of global energy inefficiency comes from the 2012 world energy balance. Of 13{,}461~Mtep of primary energy produced, only 8{,}979~Mtep reached final consumers — meaning that 4{,}482~Mtep, or roughly one-third, was lost through conversion, distribution, and internal use. The major loss contributors included power plants (2{,}486~Mtep), internal use (789~Mtep), and distribution losses (210~Mtep). This underscores the critical importance of energy efficiency alongside production capacity.

\subsection{Energy and Human Development}

Energy is not merely a technical variable — it is a fundamental driver of human quality of life. The use of energy underpins three core dimensions of human existence: human activities (movement, nutrition), human comfort (mechanisation, heating, cooling), and human health (medical care, refrigeration of medicines, hospital power). The equation \textit{energy use = quality of life} is thus not rhetorical but empirically grounded: countries with higher per capita energy consumption consistently display better human development indicators.

The five types of energy relevant to modern societies are thermal, chemical, nuclear, electrical, and mechanical (both potential and kinetic). Their valorisation depends on three key factors: the nature of the source, the efficiency of conversion, and the losses incurred along the chain. The human body itself operates at roughly 20--25\% energy conversion efficiency, analogous to a 100~W radiator.

\section{The African Energy Case}

\subsection{Electricity Access and the Darkness of Sub-Saharan Africa}

Satellite imagery of the Earth at night starkly reveals Africa's energy deficit. Whilst Europe, East Asia, and North America glow with artificial light, the African continent — particularly sub-Saharan Africa — remains largely dark. This visual contrast directly encodes a development gap.

According to IEA data (2018), electricity access rates vary dramatically across the continent. West Africa has seen relative progress, while Central Africa and East Africa lag far behind. Countries such as Nigeria (178 million without clean cooking access), Ethiopia (100 million), and the Democratic Republic of Congo (81 million) head the list of those most deprived. In terms of electricity access progress between 2000 and 2018, West Africa has made the most gains, while Central Africa has progressed the least.

\subsection{Africa's Energy Mix and Projections}

Africa's primary energy demand is dominated by biomass and bioenergy, reflecting widespread reliance on wood fuel for cooking. In sub-Saharan Africa, 65\% of the population relied on charcoal (wood charcoal, \textit{charbon de bois}) for cooking in 2018. By 2030, under announced-policies scenarios, this proportion is projected to decrease to 50\%, while under the Africa Case (an accelerated development scenario), cleaner alternatives such as LPG (liquefied petroleum gas), improved cookstoves, and other clean solutions could cover 63\% of cooking needs.

Electricity generation by fuel type also varies significantly by subregion. North Africa relies more heavily on natural gas, South Africa on coal, while sub-Saharan Africa has significant hydropower capacity. Renewables (solar PV, wind) are projected to grow substantially to 2040 under both policy scenarios, though biomass remains structurally dominant in the near term.

The IEA estimates cumulative investment needs for the Africa Case between 2019 and 2040 to exceed 3{,}000~billion USD (2018 values), covering fossil fuel production, energy supply, renewable energy, network infrastructure, and access to electricity and clean cooking.

\subsection{Clean Cooking: A Persistent Challenge}

One of Africa's most urgent energy challenges is access to clean cooking. The traditional \textit{local brazier} (\textit{braséro local}) used across central and western Africa requires 7 to 8~kg of wood to produce 1~kg of charcoal, representing a highly inefficient and deforestation-intensive process. Improved cookstove models, such as those deployed in the DRC, offer substantially better thermal efficiency.

The map of populations without access to clean cooking (IEA, 2018) shows that the overwhelming majority of affected people live in countries where more than 75\% of the population lacks clean cooking access. These include most of sub-Saharan Africa, where the link between cooking fuel and deforestation is direct and severe.

\section{Potentials and Challenges of the Democratic Republic of Congo}

\subsection{DRC: A Country of Contrasts}

The Democratic Republic of Congo (DRC) epitomises what the presentation calls the \textit{Congolese Contrast}: a country endowed with extraordinary natural resources yet plagued by deep energy poverty. Key geographic facts:

\begin{itemize}[noitemsep]
    \item Area: 2{,}354{,}310~km$^2$ (roughly equivalent to the surface area of Western Europe)
    \item 2nd largest country in Africa after Algeria
    \item Population: approximately 92{,}377{,}986 inhabitants (39 inhabitants per km$^2$)
    \item Holds 60\% of the Congo Basin forest reserve
    \item Possesses 52\% of Africa's freshwater resources
\end{itemize}

The Congo River runs 4{,}700~km — equivalent to 15 Paris-Brussels distances — and flows at 41{,}000~m$^3$/s, making it 18 times the flow of Niagara Falls and 115 times that of the Meuse River. This makes it the second largest river in the world by discharge volume, after the Amazon.

\subsection{Energy Resource Inventory}

The DRC's primary energy resources are vast and diverse, as summarised in the following table:

\begin{center}
\begin{tabular}{clll}
\toprule
\textbf{No.} & \textbf{Primary Resource} & \textbf{Potential} & \textbf{Secondary Output} \\
\midrule
1 & Biomass         & 125 million hectares     & Electricity \\
2 & Oil             & 20 billion barrels        & Fuels, Electricity \\
3 & Natural gas     & 300~Gm$^3$               & Electricity \\
4 & Solar           & 70~GW (3{,}500--6{,}750~Wh/m$^2$/day) & Electricity \\
5 & Coal            & 720~Gtonnes              & Electricity \\
6 & Uranium         & Not assessed             & Electricity \\
7 & Hydropower      & 100~GW (44~GW at Inga)   & Electricity \\
\bottomrule
\end{tabular}
\end{center}

In addition to these, the DRC holds deposits of coal, uranium, methane (dissolved in Lake Kivu), hydrocarbons, peat, oil shale, and tar sands. Yet despite this wealth, the electricity access rate stands at a maximum of 15\%, and approximately 85\% of all energy consumed is wood energy.

\subsection{Hydroelectric Potential}

The DRC's hydroelectric potential of 100~GW is the largest of any country in Africa and among the highest in the world. Of this, 44~GW is concentrated at the Inga site on the Congo River — a site equivalent in power to 40 nuclear power plants, or twice the capacity of China's Three Gorges Dam. The Inga project envisions up to eight dams (Inga 1 through 8), of which Inga 1 and 2 are already operational, and Inga 3 (11{,}050~MW) is planned.

Other significant sites include Zongo, Ruzizi, Nseke, Nzilo, and Mwadingusha. The Busanga hydroelectric power plant (240~MW) was commissioned in 2023, and the Kakobola plant (10.5~MW) is expected online in 2026.

In 2019, hydroelectricity accounted for 99.4\% of DRC's electricity production, with solar contributing 0.3\% and biomass 0.2\%, out of a total of 11.1~TWh produced. By 2024, total installed capacity reached approximately 4{,}076~MW (hydroelectric: 3{,}646.5~MW; thermal: 421.4~MW; photovoltaic: 8.7~MW), with annual electricity production rising from 12{,}460~GWh in 2020 to 13{,}625~GWh in 2024.

\subsection{Solar and Other Renewable Potentials}

Solar irradiation across the DRC ranges from 3.5 to over 5.5~kWh/m$^2$/day, with the highest values in Katanga and the Kasai regions. The total solar potential is estimated at 70~GW, offering significant decentralised generation opportunities.

The DRC also holds biogas potential (particularly in Katanga and the Kivus), agricultural residue energy potential distributed across multiple provinces, and geothermal resources along the East African Rift system (primarily in the Kivu provinces). These distributed sources are particularly relevant for rural electrification where grid extension is not economically viable.

\subsection{Energy Balance and Provincial Disparities}

The DRC's energy balance (UNDP/SIE-2011) reveals that biomass accounts for 93.3\% of total primary energy consumption, hydrocarbons for 4.1\%, and electricity for only 2.6\%. This confirms that modern energy services remain inaccessible to the vast majority of the population. Wood energy for cooking dominates household energy use nationwide.

Furthermore, electricity access is geographically highly concentrated. Data from the SIE (2014) show that Kinshasa alone accounts for around 44\% of national electricity consumption, while provinces such as Bandundu, Kasai Occidental, and Equateur each consume less than 2\%. This disparity reflects both the absence of a national grid and the near-total lack of electrification outside major urban centres.

According to ARE (2024) data, the geographic coverage rate of electricity infrastructure stands at approximately 35.9\%, while the actual electrification rate (households with electricity) is only around 10.3\% — and the true access rate (regular, reliable supply) at around 10.3\%. Electricity access has grown very slowly between 2020 and 2024.

\section{Energy Policy}

\subsection{Legal and Institutional Framework}

The DRC undertook significant energy sector reform through Law No.\ 14/011 of 17 June 2014, which liberalised the electricity sector. This law broke the monopoly previously held by SNEL (Société Nationale d'Électricité), opening the market to private and public-private partnerships. An independent regulatory authority (ARE — Autorité de Régulation de l'Électricité) was established to oversee the sector. SNEL itself has undergone reform efforts to improve its operational capacity.

\subsection{Key Policy Priorities}

The DRC's energy policy is structured around five strategic axes:

\begin{enumerate}[noitemsep]
    \item Development of the Inga hydroelectric complex (Inga 1--8, up to 44~GW), alongside smaller sites such as Kakobola, Katende, and others;
    \item Reduction of greenhouse gas emissions by 17\% through the deployment of renewable energy;
    \item Reforestation: planting 3 million forest hectares by 2030 to offset deforestation driven by charcoal production;
    \item Rural electrification through the ANSER agency (Agence Nationale de l'Électrification et des Services Énergétiques en milieu Rural et Péri-urbain), using mini-grids and individual systems;
    \item Promotion of improved cookstoves and cleaner cooking fuels to reduce wood consumption and indoor air pollution.
\end{enumerate}

\subsection{Renewable Energy and CO\texorpdfstring{\textsubscript{2}}{2} Avoidance}

The DRC's near-total reliance on hydropower for electricity generation already provides a significant carbon benefit. ARE data (2024) indicate that between 2020 and 2024, the use of renewable energy (primarily hydroelectric) avoided approximately 8.5 to 9.2~MtCO$_2$ per year, equivalent to the removal of over 200{,}000 cars from circulation annually. The number of trees that would need to be planted to achieve equivalent carbon sequestration rose from approximately 338~million in 2020 to 366~million in 2024 — underscoring the scale of the avoided emissions.

\section{Outlook and Conclusion}

\subsection{The Energy Consumption Gap}

A comparison of per capita energy consumption highlights the DRC's extreme energy poverty. Whereas Belgium and Germany consume roughly 55{,}000--80{,}000 equivalent kWh per capita per year, and South Africa around 3{,}486~kWh of electrical energy per capita, the DRC's per capita electrical energy consumption stands at approximately 374~kWh/year — less than one-ninth of South Africa's and a tiny fraction of European levels.

Two scenarios illustrate the scale of effort required:
\begin{itemize}[noitemsep]
    \item \textbf{DRC = South Africa}: matching South Africa's per capita electrical consumption would require multiplying the DRC's current output by a factor of 9.32;
    \item \textbf{DRC = Algeria}: matching Algeria's per capita consumption (1{,}267~kWh/year, largely gas-based) would require a factor of 3.39.
\end{itemize}
In both scenarios, the DRC's production would need to rely predominantly on hydropower and solar, given the country's renewable resource endowment.

\subsection{Strategic Outlook}

The DRC's energy future involves navigating a fundamental tension between \textbf{energy growth} — rapidly expanding access to meet development needs — and \textbf{energy transition} — ensuring that this growth is clean, efficient, and sustainable. The key strategic directions proposed are:

\begin{enumerate}[noitemsep]
    \item Promoting clean energy growth through smart and green technologies;
    \item Developing a diversified energy mix (hydropower, solar, gas, geothermal, biogas);
    \item Encouraging local production-consumption systems (mini-grids) where national grid extension is not feasible;
    \item Pursuing the Grand Inga Project and green hydrogen production at the Banana site (Fortescue group project);
    \item Expanding rural electrification via ANSER, including Vortex hydropower turbines for small watercourses;
    \item Pursuing the equation: \textbf{Energy = GDP = SDG = Better Quality of Life}.
\end{enumerate}

\subsection{Conclusion}

The DRC encapsulates, in extreme form, the central paradox of African energy: a continent rich in energy resources yet profoundly energy-poor. The DRC holds hydroelectric potential equivalent to 40 nuclear power plants, solar resources comparable to the world's sunniest regions, and vast biomass reserves — yet barely 10\% of its population has reliable access to electricity, and 85\% of its energy comes from wood. Globally, fossil fuels remain dominant despite three decades of renewable growth, and Africa produces only 3\% of global CO$_2$ emissions while hosting one-fifth of the world's population.

Resolving this paradox requires not only investment in infrastructure but also governance reform, technology transfer, regional integration, and a policy framework that links energy access directly to human development and climate objectives.